\section{Introduction and summary of results}
Consider a Hamburger moment sequence $(s_n)$ given as
\begin{gather*}
s_n=\int_{-\infty}^\infty x^n{\rm d}\mu(x), \qquad n\ge 0,
\end{gather*}
 where $\mu$ is a probability measure with infinite support and moments of any order. The corresponding orthonormal polynomials $P_n$, obtained from the sequence of monomials $x^n$ by the Gram--Schmidt procedure, satisfy $P_0=1$ and
\begin{gather}\label{eq:rec}
xP_n(x)=b_nP_{n+1}(x)+a_nP_n(x)+b_{n-1}P_{n-1}(x),\qquad n\ge 0,
\end{gather}
where $b_n>0, $ $a_n\in\mathbb{R}$ with the convention $b_{-1}=1$, $P_{-1}=0$. For $n\ge 0$ we have $\deg P_n=n$ and the leading coefficient of $P_n$ is equal to
\begin{gather}\label{eq:lead}
{1\over b_0b_1\cdots b_{n-1}}.
\end{gather}

We are interested in the so-called indeterminate case, when the moments $s_n$ do not determine the measure $\mu$ uniquely. Such measures must have unbounded support. In particular the numbers $s_{2n}$ must grow faster than any exponential.

The infinite Hankel matrix $\mathcal H=\{h_{m,n}\}$
 \begin{gather*}
h_{m,n}=s_{m+n}=\big\langle x^m,x^n\big\rangle_{L^2(\mu)},\qquad m,n\ge 0
\end{gather*}
is positive definite in the sense that
\begin{gather}\label{eq:Hank}
(\mathcal Hv,v)=\sum_{m,n=0}^\infty s_{m+n}v_n\overline{v}_m>0,\qquad v\in \mathcal{F}_c(\mathbb{N}_0), \qquad v\neq 0,
\end{gather}
where $\mathcal{F}_c(\mathbb{N}_0)$ denotes the set of complex sequences $v=(v_n)_{n\ge 0}$, with only finitely many non-zero entries.

Berg, Chen and Ismail \cite{B:C:I} showed that the quadratic form associated with the matrix $\mathcal H$ is bounded below, i.e., for a suitable constant $C>0$
\begin{gather*}
(\mathcal{H}v,v)\ge C\|v\|^2,\qquad v\in \mathcal{F}_c(\mathbb{N}_0),
\end{gather*}
if and only if the moment problem associated with $(s_n)$ is indeterminate.

Hamburger obtained a related characterization of indeterminacy in terms of the quadratic forms associated with the matrices $\{s_{j+k}\}$ and $\{s_{j+k+2}\}$, see \cite[p.~70]{S:T}. For a comparison between Hamburger's result and the result of Berg--Chen--Ismail, see \cite{B:C:I}.

Yafaev \cite{Y} contains a characterization of closability of the Hankel form \eqref{eq:Hank} leading to the determinate case.

The smallest eigenvalue $\lambda_N$ of the truncated matrix $\mathcal{H}_N=\{s_{m+n}\}_{0\le m,n\le N}$ has been studied in a number of papers starting with Szeg\H{o} and Widom--Wilf, see~\cite{B:S1} and references therein.

The inverse of $\mathcal H_N$ can simply be described as the matrix
\begin{gather*}
\mathcal A^{(N)}=\{a_{j,k}(N)\}_{0\le j,k\le N},
\end{gather*}
where
\begin{gather}\label{eq:N-repker}
K_N(z,w)=\sum_{n=0}^NP_n(z)P_n(w)=\sum_{j,k=0}^N a_{j,k}(N)z^jw^k.
\end{gather}
A proof and historical details can be found in \cite[Theorem~1.1]{B:S1}.

In the indeterminate case it is possible to let $N\to \infty$ in \eqref{eq:N-repker} and obtain the infinite matrix $\mathcal A=\{a_{j,k}\}$ with $a_{j,k}=\lim\limits_{N\to\infty}a_{j,k}(N)$ as coefficient matrix of the entire function of two complex variables
\begin{gather*}
K(z,w)=\sum_{n=0}^\infty P_n(z)P_n(w)=\sum_{j,k=0}^\infty a_{j,k}z^jw^k,
\end{gather*}
called the reproducing kernel of the indeterminate moment problem.

The coefficients of $\mathcal H$ can grow fast, but on the other hand the coefficients of $\mathcal A$ decay fast, so it is natural to examine if it is possible to let $N$ tend to infinity in the equations
\begin{gather*}
\sum_{k=0}^N a_{i,k}(N)s_{k+j}=\delta_{i,j},\qquad 0\le i,j\le N,
\end{gather*}
and to obtain the infinite matrix equation
\begin{gather}\label{eq:AH}
\mathcal A\mathcal H=\mathcal H\mathcal A=\mathcal I,
\end{gather}
i.e.,
\begin{gather}\label{eq:AH1}
\sum_{k=0}^\infty a_{i,k}s_{k+j} =\delta_{i,j},\qquad i,j \ge 0,
\end{gather}
meaning that all the above series are convergent with the given sum.
Notice that since both matrices $\mathcal H$ and $\mathcal A$ are symmetric, it is enough to prove $\mathcal A\mathcal H=\mathcal I$ in~\eqref{eq:AH}.

Since absolute convergence is often more easy to establish than convergence, we investigate the question if the series in~\eqref{eq:AH1} are absolutely convergent for all~$i$,~$j$.

For the Stieltjes--Wigert moment sequence (sometimes called log-normal moments)
\begin{gather*}
s_n=q^{-\tfrac12 (n+1)^2},\qquad 0<q<1,
\end{gather*}
 it was shown that \eqref{eq:AH} holds and all the series in \eqref{eq:AH1} are absolutely convergent, see \cite[Section~5]{B:S1}. After this was done the same results were established but not published for several other known indeterminate moment problems, and we began to search for general theorems about~\eqref{eq:AH}. In this paper we present the general results obtained so far. At present we have no applications of equation~\eqref{eq:AH}.

If $\mathcal H$ is replaced by an infinite positive definite Toeplitz matrix $M$, the question of existence of a classical inverse of $M$ in the sense of \eqref{eq:AH} has been studied in~\cite{E:G:T}.

The sequence $c_k:=\sqrt{a_{k,k}}$ plays a crucial role in our investigations. It decays quickly to zero in the sense that $\lim\limits_{k\to\infty} k\root{k}\of{c_{k}}=0$, see~\eqref{eq:minexpt}.

Concerning the matrix $\mathcal A$, it was shown in \cite[Section~4]{B:S1} that it is of trace class, and that
\begin{gather*}
\tr\mathcal A=\frac{1}{2\pi}\int_0^{2\pi} \sum_{k=0}^\infty \big|P_k\big({\rm e}^{{\rm i}t}\big)\big|^2{\rm d}t.
\end{gather*}

In the indeterminate case the infinite Hankel matrix $\mathcal H$ does not correspond to an operator on $\ell^2$ defined on the subspace spanned by the standard orthonormal basis $(\delta_n)_{n\ge 0}$ in $\ell^2$, where $\delta_n=(\delta_{n,m})_{m\ge 0}$.

In fact, by a theorem of Carleman we necessarily have $\sum\limits_{n=0}^\infty s_{2n}^{-1/(2n)}<\infty$, hence $s_{2n}\ge 1$ for $n$ sufficiently large, and therefore
\begin{gather*}
\sum_{m=0}^\infty s_{n+m}^2=\infty\qquad \mbox{for all} \ n.
\end{gather*}

In connection with the study of \eqref{eq:AH} we introduce the following:

\begin{defn}\label{thm:defAH} An indeterminate moment problem has property
\begin{enumerate}\itemsep=0pt
\item[{\rm (c)}] if all the series in \eqref{eq:AH1} are convergent,
\item[{\rm (ac)}] if all the series in \eqref{eq:AH1} are absolutely convergent,
\item[{\rm (aci)}] if it has property {\rm (ac)} and \eqref{eq:AH} holds,
\item[{\rm (cs)}] if $\sum\limits_{n=0}^\infty c_{2n}s_{2n}<\infty$,
\item[{\rm (cs*)}] if $\sum\limits_{l=0}^\infty c_l|s_{l+m}|<\infty$ for all $m\ge 0$.
\end{enumerate}
\end{defn}
